% This file should be rename with the speaker's last name. (e.g : smith)% Write the author names in the order you wish them to appear% and underline the one who will give the talk.% Only the speaker's email address is required.\begin{abstract_online}{Using a CAS-developed random samples generator for teaching and research in probabilistic cellular automata and Statistics}{%        \underline{Gabriel Aguilera-Venegas}$^{1}$, Jos\'e L. Gal\'an-Garc\'{\i}a$^{1}$, Mar\'{\i}a \'A. Gal\'an-Garc\'{\i}a$^{1}$, Pedro Rodr\'{\i}guez-Cielos$^{1}$, Yolanda Padilla-Dom\'{\i}nguez$^{1}$, Mar\'\i a Gal\'an-Luque$^{1}$}{[gaguilera@uma.es]        }{% $^1$ University of M\'alaga, M\'alaga, Spain% \newline{}$^2$ Physics Department, Lancaster University, Lancaster, UK}Our group introduced random samples generation using a CAS, specifically Derive, in the Derive session of ACA 2009. This talk was extended in a paper published in The Derive Newsletter [1]. In the year 2017 we presented a new version focused on research of the random samples generator in STATA that was communicated in [2]. Now a new version of the generator is being implemented in Python with the symbolic pack Simpy. In this case, our we are mainly focused  in education.We are now researching in extensions of the Conway's Game of Life, specifically using probabilistic cellular automata. Our long-term goal is the simulation of the growth of cancerous tissues. We have supervised a master thesis with the preliminary works in this topic. For the generation of random numbers involved in the probabilistic automata of this work, we have used the results developed  in the previous mentioned work. A summary of this work is about to be published in Advances in Computational Mathematics [3]. This work is both, an education and a research experience for the student who carried out the master thesis.Moreover, the experience of teaching the subject Statistics to engineering students using the material developed for random samples generation has turned out to be useful in the teaching and learning process. The marks obtained in this subject by the students are statistically significantly better than the marks obtained by students of the same subject in others group of engineering where this material has not been used. A brief statistical study about the situation is carried out.\textbf{Keywords} \newline{} Random samples generation, probabilistic automata, Game of Life, CAS.\textbf{References} \newline{}% Model for references to books.   [1] \textsc{J.L. Gal\'an, G. Aguilera, P. Rodr\'\i guez, Y. Padilla, M.A. Gal\'an}, \emph{Random distributions.mth: Random samples from distributions with Derive.}          The Derive Newsletter, 75, pages: 22-43, 2009. I.S.S.N.: 1990-7079.[2] \textsc{G. Aguilera, J.L. Gal\'an, M.A. Gal\'an, P. Rodr\'\i guez, R. Rodr\'\i guez}, \emph{Random samples generation with Stata from continuous and discrete distributions}   Proceedings of the 2017 Spanish Stata Users Group meeting. Madrid. October 19, 2017      %  Model for references to articles[3] \textsc{G. Aguilera-Venegas; J. L.Gal\'an-Garc\'\i a; R. Egea-Guerrero, M. \'A. Gal\'an-Garc\'\i a, P. Rodr\'\i guez-Cielos, Y. Padilla-Dom\'\i nguez, M. Gal\'an-Luque}, A probabilistic extension to Conway's Game of Life \emph{Advances in Computational Mathematics} (in press). DOI: https://doi.org/10.1007/S10444-019-09696-8.%  Model for references to proceedings. %[3] \textsc{A. Uthor}, Title. In \emph{Title of the proccedings}, names of the editors (eds.), first page--last page. Editor(s), City, year of publication. \newline{}\end{abstract_online}    